\documentclass[12pt,a4paper]{article}
\usepackage[utf8]{vietnam}
\usepackage{amsmath,amssymb}
\usepackage{tabularx}
\usepackage{tikz}
\usepackage{xcolor}
\usepackage[top=0.6cm, bottom=0.6cm, left=0.6cm, right=0.6cm]{geometry}
\usepackage[dethi]{ex_test}

\SAOPTN{widthSA=2.8,heightSA=0.62}
\setlength{\parskip}{0pt}

% Dong ke cham mo chuan viet tay
\newcommand{\dongke}[1]{%
  \begingroup%
  \setlength{\parskip}{0pt}%
  \baselineskip=18.5pt%
  \vspace{1.0mm}%
  \foreach \i in {1,...,#1}{%
    \noindent{\color{gray!60}\dotfill}\par%
  }%
  \vspace{0.6mm}%
  \endgroup%
}

% Dinh nghia o bai tap tu luan trong bang
\newcommand{\oBaiTap}[3][6]{%
  \vspace{1.0mm}%
  \noindent\textbf{#2} #3\par%
  \dongke{#1}%
}

\begin{document}
\pagestyle{empty}
% =============================================================================
% ========================= MA DE 101 - MAT TRUOC (TRANG 1/2) =============
% =============================================================================
\noindent
\begin{tabularx}{\linewidth}{|X|c|}
\hline
\textbf{PHIẾU KIỂM TRA ĐỊNH KỲ TOÁN 6 - TUẦN 09 - BUỔI 01} & \textbf{Mã đề: 101 - Trang 1/2} \\
\hline
\end{tabularx}

\vspace{1.0mm}

\noindent
\begin{tabularx}{\linewidth}{|p{3.2cm}|X|}
\hline
\centering\textbf{Điểm số} & \centering\textbf{Lời nhận xét của thầy cô} \tabularnewline
\rule{0pt}{1.8cm} & \\
\hline
\end{tabularx}

\vspace{1.0mm}

\noindent\textbf{Câu 1 (0,25 điểm).} Tìm $BCNN(7, 9)$: \shortans{$63$}\par\vspace{0.25mm}
\noindent\textbf{Câu 2 (0,25 điểm).} Tìm $BCNN(36, 12)$: \shortans{$36$}\par\vspace{0.25mm}
\noindent\textbf{Câu 3 (0,25 điểm).} Cho $a = 2^3 \cdot 3$ và $b = 2^2 \cdot 3^2 \cdot 5$. Giá trị $BCNN(a, b)$ là: \shortans{$360$}\par\vspace{0.25mm}
\noindent\textbf{Câu 4 (0,25 điểm).} Số tự nhiên nhỏ nhất khác 0 chia hết cho cả 6 và 8 là: \shortans{$24$}\par\vspace{0.25mm}
\noindent\textbf{Câu 5 (0,25 điểm).} Có bao nhiêu số tự nhiên khác 0 là bội chung của 4 và 6 nhỏ hơn 50? \shortans{$4$}\par\vspace{0.25mm}
\noindent\textbf{Câu 6 (0,25 điểm).} Mẫu số chung nhỏ nhất của $\frac{5}{12}$ và $\frac{7}{18}$ là: \shortans{$36$}\par\vspace{0.25mm}
\noindent\textbf{Câu 7 (0,25 điểm).} Biết $a \cdot b = 180$ và $ƯCLN(a, b) = 6$. Khi đó $BCNN(a, b)$ là: \shortans{$30$}\par\vspace{0.25mm}
\noindent\textbf{Câu 8 (0,25 điểm).} Chu vi hình tam giác đều có cạnh dài $15\text{ cm}$ là: \shortans{$45\text{ cm}$}\par\vspace{0.25mm}
\noindent\textbf{Câu 9 (0,25 điểm).} Diện tích hình thoi có hai đường chéo $12\text{ cm}$ và $15\text{ cm}$ là: \shortans{$90\text{ cm}^2$}\par\vspace{0.25mm}
\noindent\textbf{Câu 10 (0,25 điểm).} An cứ 6 ngày đi bơi, Bình cứ 8 ngày đi bơi. Sau ít nhất bao nhiêu ngày hai bạn lại cùng đi bơi? \shortans{$24\text{ ngày}$}\par\vspace{0.25mm}

\vspace{0.3mm}
\noindent\textbf{Câu 11 (3,0 điểm). Thực hiện phép tính (hợp lý nếu có thể):}
\vspace{0.4mm}

\noindent
\begin{tabularx}{\linewidth}{|X|X|}
\hline
\oBaiTap[7]{a)}{$38 + 147 + 62 + 53$} & \oBaiTap[7]{b)}{$37 \cdot 64 + 37 \cdot 36$} \\
\hline
\oBaiTap[7]{c)}{$5 \cdot 2^3 - 18 : 3^2 + 2026^0$} & \oBaiTap[7]{d)}{$160 : (25 - 9) + 4^2 \cdot 3$} \\
\hline
\oBaiTap[7]{e)}{$150 : [35 - (11 - 6)^2]$} & \oBaiTap[7]{f)}{$120 : \{50 - [35 - (10 - 5)^2]\}$} \\
\hline
\end{tabularx}

\newpage
% =============================================================================
% ========================= MA DE 101 - MAT SAU (TRANG 2/2) ==============
% =============================================================================
\noindent
\begin{tabularx}{\linewidth}{|X|}
\hline
\centering\textbf{TOÁN 6 - TUẦN 09 - BUỔI 01 | Mã đề: 101 - Trang 2/2} \\
\hline
\end{tabularx}

\vspace{1.2mm}

\noindent\textbf{Câu 12 (2,0 điểm). Tìm số tự nhiên $x$, biết:}
\vspace{0.5mm}

\noindent
\begin{tabularx}{\linewidth}{|X|X|}
\hline
\oBaiTap[5]{a)}{$3x - 14 = 31$} & \oBaiTap[5]{b)}{$4(x + 5) = 48$} \\
\hline
\oBaiTap[7]{c)}{$x \ \vdots \ 12, x \ \vdots \ 18$ và $x$ là số tự nhiên nhỏ nhất khác $0$} & \oBaiTap[7]{d)}{$x \ \vdots \ 15, x \ \vdots \ 20$ và $100 < x < 150$} \\
\hline
\end{tabularx}

\vspace{1.2mm}

\noindent
\begin{tabularx}{\linewidth}{|X|}
\hline
\vspace{0.5mm}
\textbf{Câu 13 (2,0 điểm).} Học sinh khối 6 của một trường THCS khi xếp hàng 12, hàng 15, hàng 18 đều vừa đủ hàng. Biết số học sinh trong khoảng từ 300 đến 400 học sinh.\par
\textbf{a) (1,0 điểm)} (1,0 điểm) Tìm $BCNN(12, 15, 18)$.\par
\textbf{b) (1,0 điểm)} (1,0 điểm) Tính số học sinh khối 6 của trường đó.\par
\vspace{0.4mm}
\noindent\textit{Bài giải:}\dongke{11}
\\ \hline
\end{tabularx}

\vspace{1.2mm}

\noindent
\begin{tabularx}{\linewidth}{|X|}
\hline
\vspace{0.5mm}
\textbf{Câu 14 (0,5 điểm).} Chứng minh rằng với mọi số tự nhiên $n$, hai số $2n + 1$ và $3n + 1$ là hai số nguyên tố cùng nhau.\par
\vspace{0.4mm}
\noindent\textit{Bài giải:}\dongke{7}
\\ \hline
\end{tabularx}

\newpage
% =============================================================================
% ========================= MA DE 102 - MAT TRUOC (TRANG 1/2) =============
% =============================================================================
\noindent
\begin{tabularx}{\linewidth}{|X|c|}
\hline
\textbf{PHIẾU KIỂM TRA ĐỊNH KỲ TOÁN 6 - TUẦN 09 - BUỔI 01} & \textbf{Mã đề: 102 - Trang 1/2} \\
\hline
\end{tabularx}

\vspace{1.0mm}

\noindent
\begin{tabularx}{\linewidth}{|p{3.2cm}|X|}
\hline
\centering\textbf{Điểm số} & \centering\textbf{Lời nhận xét của thầy cô} \tabularnewline
\rule{0pt}{1.8cm} & \\
\hline
\end{tabularx}

\vspace{1.0mm}

\noindent\textbf{Câu 1 (0,25 điểm).} Tìm $BCNN(8, 9)$: \shortans{$72$}\par\vspace{0.25mm}
\noindent\textbf{Câu 2 (0,25 điểm).} Tìm $BCNN(48, 16)$: \shortans{$48$}\par\vspace{0.25mm}
\noindent\textbf{Câu 3 (0,25 điểm).} Cho $a = 2^2 \cdot 5$ và $b = 2^3 \cdot 3 \cdot 5$. Giá trị $BCNN(a, b)$ là: \shortans{$120$}\par\vspace{0.25mm}
\noindent\textbf{Câu 4 (0,25 điểm).} Số tự nhiên nhỏ nhất khác 0 chia hết cho cả 8 và 10 là: \shortans{$40$}\par\vspace{0.25mm}
\noindent\textbf{Câu 5 (0,25 điểm).} Có bao nhiêu số tự nhiên khác 0 là bội chung của 6 và 8 nhỏ hơn 80? \shortans{$3$}\par\vspace{0.25mm}
\noindent\textbf{Câu 6 (0,25 điểm).} Mẫu số chung nhỏ nhất của $\frac{3}{8}$ và $\frac{5}{12}$ là: \shortans{$24$}\par\vspace{0.25mm}
\noindent\textbf{Câu 7 (0,25 điểm).} Biết $a \cdot b = 240$ và $ƯCLN(a, b) = 4$. Khi đó $BCNN(a, b)$ là: \shortans{$60$}\par\vspace{0.25mm}
\noindent\textbf{Câu 8 (0,25 điểm).} Chu vi hình tam giác đều có cạnh dài $18\text{ cm}$ là: \shortans{$54\text{ cm}$}\par\vspace{0.25mm}
\noindent\textbf{Câu 9 (0,25 điểm).} Diện tích hình thoi có hai đường chéo $14\text{ cm}$ và $10\text{ cm}$ là: \shortans{$70\text{ cm}^2$}\par\vspace{0.25mm}
\noindent\textbf{Câu 10 (0,25 điểm).} Hoa cứ 4 ngày tưới hoa, Mai cứ 6 ngày tưới hoa. Sau ít nhất bao nhiêu ngày hai bạn lại cùng tưới? \shortans{$12\text{ ngày}$}\par\vspace{0.25mm}

\vspace{0.3mm}
\noindent\textbf{Câu 11 (3,0 điểm). Thực hiện phép tính (hợp lý nếu có thể):}
\vspace{0.4mm}

\noindent
\begin{tabularx}{\linewidth}{|X|X|}
\hline
\oBaiTap[7]{a)}{$45 + 138 + 55 + 62$} & \oBaiTap[7]{b)}{$43 \cdot 72 + 43 \cdot 28$} \\
\hline
\oBaiTap[7]{c)}{$6 \cdot 2^3 - 27 : 3^2 + 2025^0$} & \oBaiTap[7]{d)}{$180 : (27 - 9) + 5^2 \cdot 2$} \\
\hline
\oBaiTap[7]{e)}{$180 : [45 - (12 - 7)^2]$} & \oBaiTap[7]{f)}{$160 : \{50 - [35 - (12 - 7)^2]\}$} \\
\hline
\end{tabularx}

\newpage
% =============================================================================
% ========================= MA DE 102 - MAT SAU (TRANG 2/2) ==============
% =============================================================================
\noindent
\begin{tabularx}{\linewidth}{|X|}
\hline
\centering\textbf{TOÁN 6 - TUẦN 09 - BUỔI 01 | Mã đề: 102 - Trang 2/2} \\
\hline
\end{tabularx}

\vspace{1.2mm}

\noindent\textbf{Câu 12 (2,0 điểm). Tìm số tự nhiên $x$, biết:}
\vspace{0.5mm}

\noindent
\begin{tabularx}{\linewidth}{|X|X|}
\hline
\oBaiTap[5]{a)}{$4x - 17 = 35$} & \oBaiTap[5]{b)}{$5(x + 6) = 65$} \\
\hline
\oBaiTap[7]{c)}{$x \ \vdots \ 15, x \ \vdots \ 20$ và $x$ là số tự nhiên nhỏ nhất khác $0$} & \oBaiTap[7]{d)}{$x \ \vdots \ 12, x \ \vdots \ 15$ và $150 < x < 200$} \\
\hline
\end{tabularx}

\vspace{1.2mm}

\noindent
\begin{tabularx}{\linewidth}{|X|}
\hline
\vspace{0.5mm}
\textbf{Câu 13 (2,0 điểm).} Học sinh khối 6 của một trường THCS khi xếp hàng 10, hàng 12, hàng 15 đều vừa đủ hàng. Biết số học sinh trong khoảng từ 200 đến 300 học sinh.\par
\textbf{a) (1,0 điểm)} (1,0 điểm) Tìm $BCNN(10, 12, 15)$.\par
\textbf{b) (1,0 điểm)} (1,0 điểm) Tính số học sinh khối 6 của trường đó.\par
\vspace{0.4mm}
\noindent\textit{Bài giải:}\dongke{11}
\\ \hline
\end{tabularx}

\vspace{1.2mm}

\noindent
\begin{tabularx}{\linewidth}{|X|}
\hline
\vspace{0.5mm}
\textbf{Câu 14 (0,5 điểm).} Chứng minh rằng với mọi số tự nhiên $n$, hai số $3n + 2$ và $4n + 3$ là hai số nguyên tố cùng nhau.\par
\vspace{0.4mm}
\noindent\textit{Bài giải:}\dongke{7}
\\ \hline
\end{tabularx}

\newpage
% =============================================================================
% ========================= MA DE 103 - MAT TRUOC (TRANG 1/2) =============
% =============================================================================
\noindent
\begin{tabularx}{\linewidth}{|X|c|}
\hline
\textbf{PHIẾU KIỂM TRA ĐỊNH KỲ TOÁN 6 - TUẦN 09 - BUỔI 01} & \textbf{Mã đề: 103 - Trang 1/2} \\
\hline
\end{tabularx}

\vspace{1.0mm}

\noindent
\begin{tabularx}{\linewidth}{|p{3.2cm}|X|}
\hline
\centering\textbf{Điểm số} & \centering\textbf{Lời nhận xét của thầy cô} \tabularnewline
\rule{0pt}{1.8cm} & \\
\hline
\end{tabularx}

\vspace{1.0mm}

\noindent\textbf{Câu 1 (0,25 điểm).} Tìm $BCNN(7, 11)$: \shortans{$77$}\par\vspace{0.25mm}
\noindent\textbf{Câu 2 (0,25 điểm).} Tìm $BCNN(60, 15)$: \shortans{$60$}\par\vspace{0.25mm}
\noindent\textbf{Câu 3 (0,25 điểm).} Cho $a = 3^2 \cdot 5$ và $b = 2 \cdot 3 \cdot 5^2$. Giá trị $BCNN(a, b)$ là: \shortans{$450$}\par\vspace{0.25mm}
\noindent\textbf{Câu 4 (0,25 điểm).} Số tự nhiên nhỏ nhất khác 0 chia hết cho cả 9 và 12 là: \shortans{$36$}\par\vspace{0.25mm}
\noindent\textbf{Câu 5 (0,25 điểm).} Có bao nhiêu số tự nhiên khác 0 là bội chung của 5 và 6 nhỏ hơn 100? \shortans{$3$}\par\vspace{0.25mm}
\noindent\textbf{Câu 6 (0,25 điểm).} Mẫu số chung nhỏ nhất của $\frac{7}{15}$ và $\frac{9}{20}$ là: \shortans{$60$}\par\vspace{0.25mm}
\noindent\textbf{Câu 7 (0,25 điểm).} Biết $a \cdot b = 300$ và $ƯCLN(a, b) = 5$. Khi đó $BCNN(a, b)$ là: \shortans{$60$}\par\vspace{0.25mm}
\noindent\textbf{Câu 8 (0,25 điểm).} Chu vi hình tam giác đều có cạnh dài $16\text{ cm}$ là: \shortans{$48\text{ cm}$}\par\vspace{0.25mm}
\noindent\textbf{Câu 9 (0,25 điểm).} Diện tích hình thoi có hai đường chéo $16\text{ cm}$ và $12\text{ cm}$ là: \shortans{$96\text{ cm}^2$}\par\vspace{0.25mm}
\noindent\textbf{Câu 10 (0,25 điểm).} Tuyến xe A cứ 10 phút, tuyến B cứ 15 phút xuất bến. Sau ít nhất bao nhiêu phút hai xe cùng xuất bến? \shortans{$30\text{ phút}$}\par\vspace{0.25mm}

\vspace{0.3mm}
\noindent\textbf{Câu 11 (3,0 điểm). Thực hiện phép tính (hợp lý nếu có thể):}
\vspace{0.4mm}

\noindent
\begin{tabularx}{\linewidth}{|X|X|}
\hline
\oBaiTap[7]{a)}{$34 + 156 + 66 + 44$} & \oBaiTap[7]{b)}{$56 \cdot 63 + 56 \cdot 37$} \\
\hline
\oBaiTap[7]{c)}{$7 \cdot 2^3 - 36 : 3^2 + 2024^0$} & \oBaiTap[7]{d)}{$140 : (26 - 12) + 3^2 \cdot 6$} \\
\hline
\oBaiTap[7]{e)}{$210 : [55 - (13 - 8)^2]$} & \oBaiTap[7]{f)}{$200 : \{50 - [35 - (13 - 8)^2]\}$} \\
\hline
\end{tabularx}

\newpage
% =============================================================================
% ========================= MA DE 103 - MAT SAU (TRANG 2/2) ==============
% =============================================================================
\noindent
\begin{tabularx}{\linewidth}{|X|}
\hline
\centering\textbf{TOÁN 6 - TUẦN 09 - BUỔI 01 | Mã đề: 103 - Trang 2/2} \\
\hline
\end{tabularx}

\vspace{1.2mm}

\noindent\textbf{Câu 12 (2,0 điểm). Tìm số tự nhiên $x$, biết:}
\vspace{0.5mm}

\noindent
\begin{tabularx}{\linewidth}{|X|X|}
\hline
\oBaiTap[5]{a)}{$5x - 19 = 46$} & \oBaiTap[5]{b)}{$3(x + 8) = 45$} \\
\hline
\oBaiTap[7]{c)}{$x \ \vdots \ 14, x \ \vdots \ 21$ và $x$ là số tự nhiên nhỏ nhất khác $0$} & \oBaiTap[7]{d)}{$x \ \vdots \ 10, x \ \vdots \ 15$ và $100 < x < 140$} \\
\hline
\end{tabularx}

\vspace{1.2mm}

\noindent
\begin{tabularx}{\linewidth}{|X|}
\hline
\vspace{0.5mm}
\textbf{Câu 13 (2,0 điểm).} Học sinh khối 6 của một trường THCS khi xếp hàng 12, hàng 15, hàng 20 đều vừa đủ hàng. Biết số học sinh trong khoảng từ 300 đến 400 học sinh.\par
\textbf{a) (1,0 điểm)} (1,0 điểm) Tìm $BCNN(12, 15, 20)$.\par
\textbf{b) (1,0 điểm)} (1,0 điểm) Tính số học sinh khối 6 của trường đó.\par
\vspace{0.4mm}
\noindent\textit{Bài giải:}\dongke{11}
\\ \hline
\end{tabularx}

\vspace{1.2mm}

\noindent
\begin{tabularx}{\linewidth}{|X|}
\hline
\vspace{0.5mm}
\textbf{Câu 14 (0,5 điểm).} Chứng minh rằng với mọi số tự nhiên $n$, hai số $2n + 3$ và $3n + 4$ là hai số nguyên tố cùng nhau.\par
\vspace{0.4mm}
\noindent\textit{Bài giải:}\dongke{7}
\\ \hline
\end{tabularx}

\newpage
% =============================================================================
% ========================= MA DE 104 - MAT TRUOC (TRANG 1/2) =============
% =============================================================================
\noindent
\begin{tabularx}{\linewidth}{|X|c|}
\hline
\textbf{PHIẾU KIỂM TRA ĐỊNH KỲ TOÁN 6 - TUẦN 09 - BUỔI 01} & \textbf{Mã đề: 104 - Trang 1/2} \\
\hline
\end{tabularx}

\vspace{1.0mm}

\noindent
\begin{tabularx}{\linewidth}{|p{3.2cm}|X|}
\hline
\centering\textbf{Điểm số} & \centering\textbf{Lời nhận xét của thầy cô} \tabularnewline
\rule{0pt}{1.8cm} & \\
\hline
\end{tabularx}

\vspace{1.0mm}

\noindent\textbf{Câu 1 (0,25 điểm).} Tìm $BCNN(5, 13)$: \shortans{$65$}\par\vspace{0.25mm}
\noindent\textbf{Câu 2 (0,25 điểm).} Tìm $BCNN(72, 18)$: \shortans{$72$}\par\vspace{0.25mm}
\noindent\textbf{Câu 3 (0,25 điểm).} Cho $a = 2 \cdot 3^2$ và $b = 2^2 \cdot 3 \cdot 7$. Giá trị $BCNN(a, b)$ là: \shortans{$252$}\par\vspace{0.25mm}
\noindent\textbf{Câu 4 (0,25 điểm).} Số tự nhiên nhỏ nhất khác 0 chia hết cho cả 6 và 9 là: \shortans{$18$}\par\vspace{0.25mm}
\noindent\textbf{Câu 5 (0,25 điểm).} Có bao nhiêu số tự nhiên khác 0 là bội chung của 4 và 10 nhỏ hơn 90? \shortans{$4$}\par\vspace{0.25mm}
\noindent\textbf{Câu 6 (0,25 điểm).} Mẫu số chung nhỏ nhất của $\frac{5}{14}$ và $\frac{3}{21}$ là: \shortans{$42$}\par\vspace{0.25mm}
\noindent\textbf{Câu 7 (0,25 điểm).} Biết $a \cdot b = 360$ và $ƯCLN(a, b) = 6$. Khi đó $BCNN(a, b)$ là: \shortans{$60$}\par\vspace{0.25mm}
\noindent\textbf{Câu 8 (0,25 điểm).} Chu vi hình tam giác đều có cạnh dài $20\text{ cm}$ là: \shortans{$60\text{ cm}$}\par\vspace{0.25mm}
\noindent\textbf{Câu 9 (0,25 điểm).} Diện tích hình thoi có hai đường chéo $18\text{ cm}$ và $10\text{ cm}$ là: \shortans{$90\text{ cm}^2$}\par\vspace{0.25mm}
\noindent\textbf{Câu 10 (0,25 điểm).} Đèn A cứ 12 giây, đèn B cứ 16 giây đổi màu. Sau ít nhất bao nhiêu giây hai đèn cùng đổi màu? \shortans{$48\text{ giây}$}\par\vspace{0.25mm}

\vspace{0.3mm}
\noindent\textbf{Câu 11 (3,0 điểm). Thực hiện phép tính (hợp lý nếu có thể):}
\vspace{0.4mm}

\noindent
\begin{tabularx}{\linewidth}{|X|X|}
\hline
\oBaiTap[7]{a)}{$47 + 128 + 53 + 72$} & \oBaiTap[7]{b)}{$48 \cdot 54 + 48 \cdot 46$} \\
\hline
\oBaiTap[7]{c)}{$8 \cdot 2^3 - 45 : 3^2 + 2023^0$} & \oBaiTap[7]{d)}{$200 : (28 - 8) + 4^2 \cdot 4$} \\
\hline
\oBaiTap[7]{e)}{$240 : [65 - (14 - 9)^2]$} & \oBaiTap[7]{f)}{$240 : \{50 - [35 - (14 - 9)^2]\}$} \\
\hline
\end{tabularx}

\newpage
% =============================================================================
% ========================= MA DE 104 - MAT SAU (TRANG 2/2) ==============
% =============================================================================
\noindent
\begin{tabularx}{\linewidth}{|X|}
\hline
\centering\textbf{TOÁN 6 - TUẦN 09 - BUỔI 01 | Mã đề: 104 - Trang 2/2} \\
\hline
\end{tabularx}

\vspace{1.2mm}

\noindent\textbf{Câu 12 (2,0 điểm). Tìm số tự nhiên $x$, biết:}
\vspace{0.5mm}

\noindent
\begin{tabularx}{\linewidth}{|X|X|}
\hline
\oBaiTap[5]{a)}{$6x - 23 = 49$} & \oBaiTap[5]{b)}{$6(x + 4) = 66$} \\
\hline
\oBaiTap[7]{c)}{$x \ \vdots \ 16, x \ \vdots \ 24$ và $x$ là số tự nhiên nhỏ nhất khác $0$} & \oBaiTap[7]{d)}{$x \ \vdots \ 12, x \ \vdots \ 18$ và $50 < x < 100$} \\
\hline
\end{tabularx}

\vspace{1.2mm}

\noindent
\begin{tabularx}{\linewidth}{|X|}
\hline
\vspace{0.5mm}
\textbf{Câu 13 (2,0 điểm).} Học sinh khối 6 của một trường THCS khi xếp hàng 15, hàng 18, hàng 20 đều vừa đủ hàng. Biết số học sinh trong khoảng từ 300 đến 400 học sinh.\par
\textbf{a) (1,0 điểm)} (1,0 điểm) Tìm $BCNN(15, 18, 20)$.\par
\textbf{b) (1,0 điểm)} (1,0 điểm) Tính số học sinh khối 6 của trường đó.\par
\vspace{0.4mm}
\noindent\textit{Bài giải:}\dongke{11}
\\ \hline
\end{tabularx}

\vspace{1.2mm}

\noindent
\begin{tabularx}{\linewidth}{|X|}
\hline
\vspace{0.5mm}
\textbf{Câu 14 (0,5 điểm).} Chứng minh rằng với mọi số tự nhiên $n$, hai số $4n + 3$ và $5n + 4$ là hai số nguyên tố cùng nhau.\par
\vspace{0.4mm}
\noindent\textit{Bài giải:}\dongke{7}
\\ \hline
\end{tabularx}
\end{document}
